\subsection{Pretend Polynomials} 
Let's review our moves for solving equations.
\begin{itemize}
\item Basic moves (add, subtract, multiply, divide, etc)
\begin{itemize}
\item Add, subtract, multiply, divide, etc.
\item If we accidentally multiply by zero, we might introduce an extraneous solution.
\end{itemize}
\item Is the variable inside a function?
\begin{itemize}
\item Remove it with an inverse.
\item Might introduce extraneous solutions, might get multiple solutions.
\end{itemize}
\item Does the variable appear more than once?
\begin{itemize}
\item Is it a polynomial?
\item Can we pretend?
\item Factor and use the zero property!
\end{itemize}
\end{itemize}
What does it look like when we ``pretend'' we have a polynomial?
\begin{example}
Solve $x^{4/3}+5x^{2/3}-36=0$
\begin{lpic}{}{6}
\end{lpic}
\end{example}
\begin{example}
Solve $x^{8/3}+2x^{5/3}-8x^{2/3}=0$
\begin{lpic}{}{6}
\end{lpic}
\end{example}